\section{Two-points correlation function and power spectrum}
	Overdensity $\delta(\mathbf{x}) = (\rho(\textbf{x})-\bar{\rho})/\bar{\rho}$ describes the relative excess density at point $\mathbf{x}$.Based on the overdensity $\delta(\textbf{x})$, we obtain the two-points correlation function $\xi(r)$:
	\begin{equation}
		\xi(\abs{\mathbf{x_1}-\mathbf{x_2}}) 
		=
		\left<
			\delta(\mathbf{x_1})\delta(\mathbf{x_2})
		\right>.
	\end{equation}
	It describes the excess probability of finding two galaxies separated by distance $r=\abs{\mathbf{x_1}-\mathbf{x_2}}$. Here, 'excess' refers to the additional probability relative to a random distribution. Mathematically, the probability $\dd{P_{excess}} = n^2\xi(r)\dd{V_1}\dd{V_2}$ is the excess probability in the actual probability $\dd{P_{actual}} = n^2(1+\xi(r))\dd{V_1}\dd{V_2}$.
	
	For the ionization field, it can be written as:
	\[
		x(\mathbf{r}) = x_e + \delta x(\mathbf{r}).
	\]
	so that,
	\begin{align}
		\delta_{\text{HII}}(\mathbf{r}) &= x(\mathbf{r})\times\dfrac{\rho(\mathbf{r})}{\bar{\rho}}-x_e*\dfrac{\bar{\rho}}{\bar{\rho}}
		\notag\\
		&= 
		[x_e + \delta x(\mathbf{r})]\times[1+\delta(\mathbf{r})]-x_e
		\notag\\
		&=
		x_e + \delta x(\mathbf{r}) + [x_e + \delta x(\mathbf{r})]\delta(\mathbf{r})-x_e
		\notag\\
		&=
		\delta x(\mathbf{r}) + [x_e + \delta x(\mathbf{r})]\delta(\mathbf{r})
	\end{align}
	
	Therefore, the two-points correlation function of ionized hydrogen is,
	\begin{align}
		\xi_{\text{HII}\text{HII}}(r) &= \left<
			\delta_{\text{HII}}(\mathbf{r_1})\delta_{\text{HII}}(\mathbf{r_2})
		\right>
		\notag\\
		&=
		\left<
			\left\lbrace
			\delta x(\mathbf{r_1}) + [x_e + \delta x(\mathbf{r_1})]\delta(\mathbf{r_1})
			\right\rbrace\left\lbrace
			\delta x(\mathbf{r_2}) + [x_e + \delta x(\mathbf{r_2})]\delta(\mathbf{r_2})
			\right\rbrace
		\right>
		\notag\\
		&=
		\left<
			\delta x(\mathbf{r_1})\delta x(\mathbf{r_2})
		\right>
		+x_e^2
		\left<
			\delta(\mathbf{r_1})\delta(\mathbf{r_2})
		\right>
		+\left<
			\delta x(\mathbf{r_1})\delta(\mathbf{r_1})\delta x(\mathbf{r_2})\delta(\mathbf{r_2})
		\right>
		\notag\\
		&+
		2\left<
			\delta x(\mathbf{r_1})\delta(\mathbf{r_1})\delta x(\mathbf{r_2})
		\right>
		+2x_e\left<
		\delta x(\mathbf{r_1})\delta(\mathbf{r_1})\delta (\mathbf{r_2})
		\right>
		+2x_e\left<
		\delta(\mathbf{r_1})\delta x(\mathbf{r_2})
		\right>
	\end{align}

Considering Wick's theorem, we can decompose the four-point correlation function into two-point correlation functions. Therefore, we have:
	\begin{align}
		\xi_{\text{HII}\text{HII}}(r) &= \left<
			\delta_{\text{HII}}(\mathbf{r_1})\delta_{\text{HII}}(\mathbf{r_2})
		\right>
		\notag\\
		&=
		\left<
			\delta x(\mathbf{r_1})\delta x(\mathbf{r_2})
		\right>
		+x_e^2
		\left<
			\delta(\mathbf{r_1})\delta(\mathbf{r_2})
		\right>
		+\left<
			\delta x(\mathbf{r_1})\delta(\mathbf{r_1})
		\right>\left<
			\delta x(\mathbf{r_2})\delta(\mathbf{r_2})
		\right>
		\notag\\
		&+
		2\left<
			\delta x(\mathbf{r_1})\delta(\mathbf{r_1})
		\right>\left<
			\delta x(\mathbf{r_2})
		\right>
		+2x_e\left<
		\delta x(\mathbf{r_1})\delta(\mathbf{r_1})
		\right>\left<
		\delta (\mathbf{r_2})
		\right>
		+2x_e\left<
		\delta(\mathbf{r_1})\right>\left<\delta x(\mathbf{r_2})
		\right>
	\end{align}	